\section{Introduction}
\label{sec:intro}

Vision--language models are most useful in exactly the conditions that most degrade
an image. Assistive applications answer questions about photographs taken by blind
and low-vision users, who cannot verify whether a shot is focused, framed, or
lit~\cite{gurari2018vizwiz}; perception systems in vehicles must see through rain, fog,
glare and motion; and scientific and clinical tools routinely operate on noisy scans.
What unites these settings is a single property: a confident wrong answer is the most
damaging possible output, because whoever acts on it receives no signal that it is
wrong. When an input is genuinely uninterpretable, the correct behavior is
\emph{abstention}---stating that the image does not support an answer.

VLMs are notoriously bad at this. Across a growing body of
benchmarks~\cite{miyai2024upd,zhu2025mohobench,saadat2025visiontrap,
madhusudhan2026knowing}, state-of-the-art models answer unanswerable visual questions
rather than withhold, and recent work shows the deficit is not one of perception but of
expression: a model's hidden states often encode that an input is unanswerable while its
output commits to an answer anyway. This \emph{representation--output gap} has been
documented from several directions---linear probes that decode answerability which the
response ignores~\cite{he2026knowing,pramono2026trapsbench}, calibration
metrics~\cite{chandu2025certainly}, and the broader finding that language models
internally ``know what they know'' more than they say~\cite{kadavath2022know,
burns2023discovering}. What this literature has in common is that it characterizes the
gap in \emph{off-the-shelf} models. It leaves open a question that matters the moment one
tries to deploy a VLM in degraded conditions: what does \emph{training} the model for
those conditions do to the gap?

This is the question we study, and the answer is unwelcome (\cref{fig:teaser}). The
natural way to make a VLM robust to degraded inputs is to fine-tune it on corrupted
images, and two paradigms dominate practice. Knowledge distillation trains a student
that sees a degraded image to imitate a larger teacher that sees the clean one; it is
inexpensive, needs no ground-truth labels, and uses only teacher forward passes.
Reinforcement learning with verifiable rewards, here Group Relative Policy Optimization
(GRPO)~\cite{shao2024deepseekmath}, rewards the model when its answer matches the gold
answer. Both reliably raise accuracy on corrupted inputs---and both, we find, do
something their accuracy never reveals: they teach the model to stop abstaining. The gap
that prior work measured in static models is not fixed; a routine, well-motivated
training intervention actively widens it. This is an instance of a pattern now
familiar on the language side, where fine-tuning for one objective silently erodes a
capability it never touched~\cite{qi2024finetuning,luo2023empirical}; we give the first
account of it for visual abstention, and trace it to the training objective.

We establish this on three 4B VLMs (Qwen3-VL-4B, InternVL3.5-4B and Gemma-3-4B),
fine-tuned with on-policy reverse-KL distillation (RKL)~\cite{gu2024minillm} and with
GRPO on 30k GQA questions~\cite{hudson2019gqa} corrupted by a 16-type engine. We
evaluate on corrupted TextVQA~\cite{singh2019textvqa}, GQA~\cite{hudson2019gqa} and
VQAv2~\cite{goyal2017vqav2}, following the corruption protocol of
Usama~\etal~\cite{usama2025analysing}, and on VizWiz~\cite{gurari2018vizwiz}: real
blind-user photographs, a third of whose questions are genuinely unanswerable. We
organize the analysis around three research questions:
\textbf{RQ1}: does robust fine-tuning erase the model's internal representation of
answerability, or only suppress its expression?
\textbf{RQ2}: do RKL and GRPO fail for the same reason?
\textbf{RQ3}: can abstention be recovered without giving up the robustness gains
(\cref{eq:recover})?
The pattern we uncover is not a quirk of one recipe but a coupling between two objectives
that the field optimizes and reports in isolation. We make four contributions.

% Numbers in items 1-3: Tabs. 3-5 and 7 (no results file in this checkout; add the paths).
\begin{enumerate}
\item \textbf{A silent abstention regression.} Corruption-robust fine-tuning raises
clean and corrupted accuracy---GRPO by $7.6$ corrupted points on Qwen3-VL-4B, with gains
that transfer zero-shot to unseen corruptions (\cref{sec:robustness})---yet on all three
backbones it lowers correct refusals on unanswerable VizWiz questions (Qwen3-VL-4B:
$71.6\%$ for the base model, $49.6\%$ after RKL, $44.9\%$ after GRPO;
\cref{sec:regression}). GRPO, the most accurate model, refuses least, and no accuracy
metric exposes this.

\item \textbf{The ability is suppressed, not erased (RQ1).} A linear probe decodes
answerability from frozen hidden states at $\mathrm{AUROC}\approx0.93$ in the base, RKL
and GRPO models of Qwen3-VL-4B, and still flags $76$--$77\%$ of the unanswerable inputs
that the two fine-tuned models wrongly answer; the failure is behavioral, not representational (\cref{sec:rq1}).

\item \textbf{Two opposite failure mechanisms (RQ2).} A reasoning-transplant test shows
RKL fails on the reasoning side (given a ``can't-tell'' chain it abstains $98.4\%$ of
the time, so unaided it never reaches that chain), while GRPO fails on the answer side
(given the same chain it still answers $91.8\%$ of the time, a reward-induced shortcut)
(\cref{sec:rq2}).

\item \textbf{Two mechanism-matched remedies (RQ3).} A $10\%$ unanswerable data mix and
inference-time activation steering each restore abstention without surrendering the
corruption gains, and the objective decides which is better---entropy-gated steering
for RKL, the missing data signal for GRPO (\cref{sec:rq3}).
\TODO{Add numbers once the headline remedy and the reported probe cutoff are fixed:
Tab.~4 favors the probe switch, while this item claims gated steering for RKL and the
data mix for GRPO.}
\end{enumerate}

The lesson is methodological: corruption robustness and calibrated abstention are
coupled, standard accuracy benchmarks hide the coupling, and the right fix depends on
how the model was trained. We therefore report an explicit abstention axis beside
accuracy throughout.
